\begingroup
\small
\setlength{\tabcolsep}{3pt}
\begin{longtable}{@{}L{0.13\textwidth} L{0.19\textwidth} L{0.20\textwidth} L{0.38\textwidth}@{}}
\caption{Initial conditions used by the active TRM process.}
\label{tab:initial-conditions}\\
\toprule
Variable & XML parameter & Active value & Interpretation and modelling gate\\
\midrule
\endfirsthead
\toprule
Variable & XML parameter & Active value & Interpretation and modelling gate\\
\midrule
\endhead
\bottomrule
\endlastfoot
\(\uu(\xx,0)\) & \texttt{ic\_displacement} & \((0,0)\,\mathrm{m}\) &
The model starts from zero displacement on the already open two-dimensional geometry.
No excavation/driving phase is represented explicitly.\\
\(\pl(\xx,0)\) & \texttt{ic\_pressure} &
\(1.5\times 10^{6}\,\mathrm{Pa}\) everywhere &
The XML function contains a vertical term multiplied by zero, therefore the active
initial liquid pressure is spatially uniform.  This remains a gate for pressure
history: no post-excavation drained niche field is encoded.\\
\(\TT(\xx,0)\) & \texttt{ic\_temperature} &
\(298.15\,\mathrm{K}\) &
The model is initialized at \(T_0\).  Together with the \texttt{bulk\_all}
temperature Dirichlet condition, this enforces homogeneous constant temperature in
the active parametrization.\\
\(\sig_0\) & \texttt{ic\_sigma0} &
Defined but inactive &
The file defines in-situ stress-like expressions, but the process XML comments out
the initial-stress reference.  Prestress is therefore a provenance question, not an
active initial condition in the present run.\\
\end{longtable}
\endgroup
